\documentclass[12pt,a4paper]{article}

\usepackage[spanish,es-nodecimaldot]{babel}
\usepackage[utf8]{inputenc}
\usepackage[T1]{fontenc}
\usepackage{lmodern}
\usepackage[margin=2.5cm]{geometry}
\usepackage{microtype}
\usepackage{booktabs,longtable,array,tabularx}
\usepackage{xcolor}
\usepackage{enumitem}
\usepackage{hyperref}
\usepackage{fancyhdr}
\usepackage{listings}

\hypersetup{colorlinks=true,linkcolor=blue!50!black,urlcolor=blue!50!black}
\setlist{itemsep=2pt,topsep=4pt}
\pagestyle{fancy}
\fancyhf{}
\lhead{FaroNet -- Proyecto 1}
\rhead{Documentación de avance}
\cfoot{\thepage}
\setlength{\headheight}{14.5pt}

\lstdefinestyle{terminal}{
  basicstyle=\ttfamily\small,
  breaklines=true,
  frame=single,
  columns=fullflexible,
  keepspaces=true,
  showstringspaces=false
}

\title{\textbf{FaroNet}\\Documentación del proyecto}
\author{Esteban Castro \and Esteban Espinoza\\[4pt]
Ingeniería de Computación -- Redes de Computadoras\\
Profesor: Kevin Moraga}
\date{1 de noviembre de 2026\\Versión de avance}

\begin{document}
\maketitle
\begin{center}
\fbox{\parbox{0.9\textwidth}{\centering
\textbf{Documento de avance.} El proyecto no está concluido. Este texto describe
el diseño y el estado observado al 1 de noviembre de 2026. Los resultados de
hardware, pruebas de rendimiento y despliegue remoto que no se han medido se
identifican como pendientes; no se presentan como resultados obtenidos.}}
\end{center}
\tableofcontents
\newpage

\section{Introducción}

FaroNet es un prototipo académico de comunicación que combina un enlace óptico
entre dispositivos Faro, una red mesh y servicios públicos desplegados en el
edge. El objetivo de diseño es permitir el intercambio de mensajes privados
cifrados y, de forma separada, ofrecer salas y feeds públicos. La arquitectura
busca limitar la información que conocen los nodos intermedios y evitar que el
Worker almacene mensajes privados en claro.

La solución se divide en cuatro componentes: (1) un Faro basado en Raspberry Pi
con LED emisor y receptor acondicionado; (2) software Python para framing,
acceso al medio, criptografía y transporte mesh; (3) un gateway Python que
envía operaciones autorizadas al Worker por HTTPS; y (4) un Worker JavaScript
con Durable Objects y KV. El medio óptico y el Worker se pueden desarrollar y
probar por separado antes de integrarlos.

\begin{center}
\begin{tabular}{c}
Nodo / aplicación FaroNet $\leftrightarrow$ Faro óptico $\leftrightarrow$ Faro / nodo\\
$\hspace{2cm}\updownarrow$ mesh con rutas onion de tres relays $\updownarrow$\\
Gateway local $\xrightarrow{\text{HTTPS firmado}}$ Cloudflare Worker\\
\hspace{2cm}salas públicas, feeds, dead drops y directorio
\end{tabular}
\end{center}

\section{Ambiente de desarrollo}

\subsection{Software}
\begin{itemize}
  \item Python 3.11 o posterior para Faro, nodos y gateway.
  \item PyNaCl/libsodium para cifrado autenticado, cajas selladas y firmas Ed25519.
  \item FastAPI, Uvicorn y HTTPX para el servicio local del gateway.
  \item JavaScript en Cloudflare Workers; Wrangler 4 para desarrollo local y despliegue.
  \item Durable Objects para coordinar WebSockets de salas y mantener estado temporal de autenticación y limitación.
  \item Cloudflare KV para drops, feeds, claves públicas de gateway y réplica del directorio.
  \item Windows/PowerShell para el desarrollo actual. El backend GPIO requiere Raspberry Pi OS y GPIO Zero.
\end{itemize}

\subsection{Hardware y herramientas}
El netlist preliminar utiliza Raspberry Pi con GPIO17 para TX, GPIO27 para RX,
GPIO22 para el botón de pánico y GPIO23--25 para indicadores. El circuito
incluye LED emisor con transistor, fototransistor, comparador, resistencias,
capacitores y pulsador. La placa final debe soldarse en perfboard/stripboard o
PCB y usar headers para retirar la Raspberry Pi. Se requieren cautín, estaño,
multímetro y el equipo para tomar fotografías de ambas caras de la placa.

El netlist y la lista de piezas disponibles están en
\texttt{EsquématicoCircuito/faro\_netlist\_y\_lista.md}. La documentación de
avance no afirma que la placa final esté soldada ni que se hayan completado las
mediciones eléctricas.

\subsection{Depuración y flujo de trabajo}
Se utiliza el loopback para ejercitar framing sin hardware, mensajes genéricos
de error para no exponer secretos y Wrangler local para inspeccionar las rutas
del Worker. Las claves privadas se suministran mediante variables de entorno;
no deben agregarse al repositorio. Antes de la entrega se recomienda registrar
cambios con Git, ejecutar las verificaciones de cada capa y conservar las
mediciones de hardware con fecha, distancia y condiciones ambientales.

\section{Diseño}

\subsection{Capa física y enlace}
El diseño propone OOK con codificación Manchester, transmisión half-duplex y
tramas de hasta 128 bytes antes de codificar símbolos. La trama contiene
preambulo \texttt{AA AA}, delimitador \texttt{D5}, versión, tipo, direcciones
de 8 bytes, secuencia, longitud, payload y CRC-16/CCITT-FALSE. El payload
máximo definido es 102 bytes. Los tipos iniciales son DATA, ACK, BEACON y
CONTROL. El acceso al medio realiza detección de portadora, backoff y
confirmación ACK con reintentos.

El framing y Manchester están separados de los backends físicos. Hay un
backend loopback y un backend GPIO para los pines indicados. El backend GPIO
usa temporización de software; su tasa efectiva, BER, sincronización del
receptor y desempeño frente a luz ambiental siguen pendientes de validación
en el circuito real. No se reportan como demostradas las distancias mínimas ni
las mediciones exigidas por el enunciado.

\subsection{Mesh, privacidad y aplicación privada}
Las entradas de directorio contienen clave pública, capacidades, expiración y
firma Ed25519. El nodo puede rechazar entradas inválidas o vencidas. La
búsqueda de rutas admite un mínimo configurable de relays y el transporte
privado exige al menos tres. Cada relay retira una capa onion y reenvía el
paquete interno. El contenido se cifra extremo a extremo con
XChaCha20-Poly1305 mediante PyNaCl; se usa padding de tamaño fijo para los
mensajes que caben en el límite implementado. Los mensajes privados mayores
que ese límite se rechazan.

La interfaz de aplicación privada convierte texto UTF-8 en tramas mesh y
descifra el payload recibido. La integración completa de selección de ruta,
transporte entre nodos físicos, sesiones de usuario y una interfaz de chat no
está concluida.

El controlador de pánico permite registrar buffers mutables de secretos,
cancelar estado pendiente, limpiar sesiones, rotar la dirección y activar el
indicador. El backend GPIO detecta el botón. Falta validar en hardware que la
acción física active todo el ciclo de pánico de extremo a extremo. El borrado
en Python no puede garantizar la eliminación de copias inmutables de bytes ni
de datos ya observados por terceros.

\subsection{Gateway y edge}
El gateway local expone un servicio FastAPI en loopback y firma solicitudes
HTTPS con Ed25519, timestamp y nonce. Su clave pública se registra en el
binding KV \texttt{GATEWAY\_KEYS}; la privada y el secreto compartido de drops
se configuran localmente. No se deben registrar claves, payloads, direcciones
físicas ni identificadores.

El Worker implementa una página de sala pública, salas WebSocket coordinadas
por Durable Objects, feeds Atom, directorio de lectura pública y drops
cifrados de 1024 bytes con expiración. El token del drop se deriva mediante
HMAC-SHA256 de un secreto compartido y un identificador aleatorio. El Worker
almacena el blob cifrado, sin descifrarlo. El rate limiting tiene un límite
global demostrativo. KV tiene consistencia eventual; en particular, la lectura
y eliminación de un drop no son una operación atómica frente a solicitudes
concurrentes.

\subsection{Amenazas, privacidad y límites}
Se consideran nodos mesh maliciosos, observadores del medio óptico, el operador
del gateway, el proveedor del edge y un atacante con acceso temporal al
dispositivo. Se protegen el texto privado, las claves y la relación origen--
destino. Las mitigaciones previstas son cifrado autenticado, onion routing de
tres relays, firmas verificables, expiración, padding, rotación de direcciones,
no persistencia deliberada de encabezados identificadores y botón de pánico.

Estas medidas no eliminan el análisis temporal o volumétrico del tráfico, la
observación de metadatos por el proveedor de edge, ni el compromiso de un
extremo. El Worker no escribe deliberadamente IP, User-Agent ni direcciones
físicas en almacenamiento de aplicación; el proveedor de Cloudflare puede
observar datos de conexión conforme a su plataforma. La política de no
registro debe revisarse también en las herramientas de despliegue y
observabilidad.

\section{Estructuras de datos y funciones principales}

\begin{longtable}{p{0.28\textwidth}p{0.64\textwidth}}
\toprule
\textbf{Componente} & \textbf{Estructuras y responsabilidades}\\
\midrule
\endhead
\texttt{faro.link.frame} & \texttt{Frame}, tipos, serialización, validación de longitudes y CRC.\\
\texttt{faro.link.mac} & Acceso half-duplex, CCA, backoff, ACK y cola temporal de tramas recibidas.\\
\texttt{faro.physical} & Contrato \texttt{PhysicalBackend}, loopback, Manchester y backend GPIO Raspberry Pi.\\
\texttt{faro.privacy.crypto} & Cifrado autenticado, cajas selladas y funciones de firma/verificación usando PyNaCl.\\
\texttt{faro.mesh.directory} & Entradas firmadas, canonicalización y verificación de firma y expiración.\\
\texttt{faro.mesh.routing} & Búsqueda de ruta con mínimo de relays; error explícito si la ruta no está disponible.\\
\texttt{faro.mesh.fragmentation} & Identificador de mensaje, índice/cantidad de fragmentos, reensamblado con TTL y límites de memoria.\\
\texttt{faro.mesh.onion/transport} & Capas por relay, cifrado extremo a extremo, padding y creación de tramas.\\
\texttt{faro.application.private\_chat} & Conversión texto/tramas y apertura de mensajes privados en el destino.\\
\texttt{gateway} & Cliente HTTP firmado y servicio local para feeds y drops.\\
\texttt{edge/src} & Enrutamiento Worker; salas WebSocket, rate limiter, estado anti-replay, feeds, drops y directorio.\\
\bottomrule
\end{longtable}

\section{Instrucciones de ejecución}

\subsection{Worker local}
Desde \texttt{Progra/edge}, instalar Wrangler si no está disponible y completar
los IDs de namespaces en \texttt{wrangler.toml}. En desarrollo local se
pueden crear namespaces y usar sus bindings locales. La clave pública del
gateway se coloca en el KV local \texttt{GATEWAY\_KEYS}; los valores privados
no se escriben en el repositorio.

\begin{lstlisting}[style=terminal]
npx wrangler kv namespace create DROPS
npx wrangler kv namespace create DIRECTORY
npx wrangler kv namespace create PUBLIC_FEEDS
npx wrangler kv namespace create GATEWAY_KEYS
npx wrangler dev --local-protocol https
\end{lstlisting}

La URL local reportada durante la sesión de desarrollo fue
\texttt{https://127.0.0.1:8787}. Se observó respuesta HTTP 200 para \texttt{/}
y \texttt{/v1/nodes}. Esta comprobación confirma el inicio y esas dos rutas;
no equivale a una prueba integral de autenticación, WebSockets, expiración o
despliegue remoto.

\subsection{Gateway local}
Desde \texttt{Progra/gateway}, crear/activar el entorno virtual e instalar el
proyecto:
\begin{lstlisting}[style=terminal]
py -m venv .venv
.\.venv\Scripts\Activate.ps1
python -m pip install -e .
\end{lstlisting}

En la misma terminal, configurar las variables requeridas sin guardar valores
secretos en archivos versionados:
\begin{lstlisting}[style=terminal]
$env:FARONET_EDGE_URL = "https://127.0.0.1:8787"
$env:FARONET_KEY_ID = "ID_PUBLICO_DEL_GATEWAY"
$env:FARONET_SIGNING_KEY_HEX = "CLAVE_PRIVADA_HEX"
$env:FARONET_DROP_SECRET_B64 = "SECRETO_COMPARTIDO_BASE64URL"
python -m gateway.service
\end{lstlisting}

La terminal con Wrangler debe permanecer abierta. La ejecución local del
gateway debe considerarse pendiente hasta obtener una respuesta satisfactoria
de \texttt{/health} y una publicación autenticada de prueba.

\subsection{Faro y módulos Python}
Desde \texttt{Progra/faro}, instalar el paquete. En Raspberry Pi, instalar
además el extra GPIO. El backend debe instanciarse con la tasa de símbolos
medida para el circuito y enlazarse con \texttt{OpticalLink}.
\begin{lstlisting}[style=terminal]
python -m pip install -e .
python -m pip install -e ".[rpi]"
\end{lstlisting}

La ejecución con GPIO no se considera validada hasta intercambiar tramas entre
dos Faros y registrar resultados a tres distancias con luz ambiental normal.

\section{Actividades realizadas por estudiante}

\begin{center}
\fcolorbox{red}{red!5}{\parbox{0.9\textwidth}{\textbf{Borrador ilustrativo, pendiente de validación.} La tabla siguiente es una
propuesta de organización basada en la bitácora de referencia y el estado del
repositorio. Las fechas, horas y atribuciones no son registros comprobados.
Esteban Castro y Esteban Espinoza deben corregirla usando sus actividades
reales antes de cualquier entrega académica.}}
\end{center}

\begin{longtable}{p{0.17\textwidth}p{0.15\textwidth}p{0.11\textwidth}p{0.48\textwidth}}
\toprule
\textbf{Fecha y hora} & \textbf{Estudiante} & \textbf{Horas estimadas} & \textbf{Actividad propuesta para validar}\\
\midrule
\endhead
29/10/2026, 18:30--20:00 & Esteban Castro & 1.5 & Revisión del enunciado, requisitos del Faro físico y división inicial de tareas.\\
29/10/2026, 18:30--20:30 & Esteban Espinoza & 2.0 & Revisión de arquitectura, capas de red y criterios de privacidad del proyecto.\\
30/10/2026, 15:00--18:30 & Esteban Castro & 3.5 & Revisión de materiales disponibles y compra de componentes faltantes para el prototipo.\\
30/10/2026, 19:00--20:30 & Esteban Espinoza & 1.5 & Organización del plan de trabajo y revisión de alternativas para Worker y gateway.\\
31/10/2026, 09:00--12:00 & Esteban Castro & 3.0 & Elaboración/revisión del netlist y esquema preliminar con Raspberry Pi y circuito óptico.\\
31/10/2026, 10:00--12:00 & Esteban Espinoza & 2.0 & Revisión de contratos de software y separación entre capas física, enlace y mesh.\\
31/10/2026, 12:40--14:10 & Esteban Castro & 1.5 & Redacción inicial del modelo de amenazas y discusión de direcciones rotables.\\
31/10/2026, 12:40--14:10 & Esteban Espinoza & 1.5 & Revisión de adversarios, mitigaciones y límites del modelo de amenazas.\\
31/10/2026, 14:30--18:00 & Esteban Castro & 3.5 & Implementación/revisión de framing, CRC, Manchester y acceso al medio en software.\\
31/10/2026, 14:30--18:00 & Esteban Espinoza & 3.5 & Implementación/revisión de directorio firmado, rutas, fragmentación y cifrado onion.\\
01/11/2026, 10:00--13:00 & Esteban Castro & 3.0 & Integración de privacidad y transporte mesh; revisión de padding y manejo de mensajes.\\
01/11/2026, 10:00--13:00 & Esteban Espinoza & 3.0 & Revisión de gateway, autenticación por firma y configuración local del Worker.\\
01/11/2026, 14:00--17:00 & Esteban Castro & 3.0 & Trabajo inicial de Worker, feeds/drops y backend GPIO experimental.\\
01/11/2026, 14:00--17:00 & Esteban Espinoza & 3.0 & Integración local, revisión de rutas HTTP y configuración de entorno del gateway.\\
01/11/2026, 17:30--18:30 & Esteban Castro & 1.0 & Consolidación de documentación de avance y pendientes de verificación.\\
01/11/2026, 17:30--18:30 & Esteban Espinoza & 1.0 & Revisión cruzada del estado del proyecto y preparación de próximos pasos.\\
\midrule
\multicolumn{2}{r}{Total ilustrativo Esteban Castro} & 20.0 & Debe reemplazarse por el total de horas reales.\\
\multicolumn{2}{r}{Total ilustrativo Esteban Espinoza} & 17.5 & Debe reemplazarse por el total de horas reales.\\
\bottomrule
\end{longtable}

\section{Estado, autoevaluación y limitaciones}

\subsection{Estado al 1 de noviembre de 2026}
El código cubre un prototipo de framing y acceso al medio, estructuras de
mesh, operaciones criptográficas, una interfaz privada mínima, un gateway
local y rutas del Worker. El Worker inició localmente con Durable Objects y
namespaces KV; las rutas \texttt{/} y \texttt{/v1/nodes} respondieron 200.
Wrangler indicó recursos en modo local. No se ha confirmado un despliegue
remoto. En la sesión documentada, el servicio gateway aún no había completado
un arranque satisfactorio.

No se han documentado resultados de BER o throughput, distancias medidas,
pruebas con dos dispositivos ópticos, placa soldada final, pánico físico de
extremo a extremo ni sesión de chat privada entre nodos físicos. La
limitación de tiempo de GPIO desde espacio de usuario y la consistencia
eventual de KV deben tenerse presentes al interpretar el prototipo.

\subsection{Autoevaluación preliminar}
Escala: 1 = muy malo, 2 = malo, 3 = regular, 4 = bueno, 5 = muy bueno. Estas
calificaciones son una autoevaluación de avance por completar por el grupo;
no son una evaluación externa.

\begin{center}
\begin{tabular}{p{0.72\textwidth}c}
\toprule
\textbf{Rubro} & \textbf{Borrador}\\
\midrule
Aprendizaje de soldadura y hardware & Pendiente\\
Aprendizaje de la capa física & Pendiente\\
Aprendizaje de enrutamiento en redes mesh & 3\\
Aprendizaje de privacidad y anonimato & 3\\
Aprendizaje de serverless en el edge & 3\\
Aprendizaje de redacción de RFCs & 3\\
\bottomrule
\end{tabular}
\end{center}
Las notas deben ser acordadas por ambos integrantes después de revisar los
resultados reales. El reporte de commits de Git se agregará cuando el código
esté en un repositorio Git con el historial correspondiente.

\subsection{Problemas y trabajo pendiente}
\begin{enumerate}
  \item Completar instalación y validación del gateway local, incluyendo una publicación autenticada y rechazo de replay.
  \item Crear namespaces remotos, reemplazar IDs locales en Wrangler, autenticar la cuenta Cloudflare y verificar despliegue reproducible.
  \item Corroborar salas WSS desde navegador, feeds Atom, tamaño fijo/TTL/consumo de drops y acceso al directorio.
  \item Integrar las aplicaciones con el transporte real de nodos y probar rutas de al menos tres relays.
  \item Calibrar el GPIO en el circuito final; medir tasa efectiva y BER en al menos tres distancias con luz ambiental normal.
  \item Soldar y revisar el Faro; registrar continuidad, ausencia de corto VCC--GND, voltajes clave, BOM final y fotografías de ambas caras.
  \item Verificar el botón de pánico, rotación de dirección y borrado de estado sensible con el dispositivo real.
  \item Actualizar RFC y documentación con mediciones, pruebas, comandos y limitaciones observadas.
\end{enumerate}

\section{Lecciones aprendidas}
\begin{itemize}
  \item Definir temprano interfaces entre framing, medio físico y aplicación permite desarrollar con loopback mientras el circuito se termina.
  \item La temporización de GPIO desde un sistema operativo general puede limitar la confiabilidad; el rendimiento debe medirse en el hardware final y no inferirse del simulador.
  \item Las operaciones de almacenamiento con consistencia eventual requieren límites claros: TTL mínimo, concurrencia y semántica de lectura única.
  \item Los secretos deben configurarse en el entorno de ejecución y nunca pegarse en documentos, bitácoras o repositorios.
  \item Una función que compila o responde localmente no demuestra por sí sola el cumplimiento integral; cada requisito debe enlazarse con evidencia reproducible.
\end{itemize}

\end{document}
